\documentclass[11pt]{article}
\usepackage[margin=1in]{geometry}
\usepackage{amsmath,amssymb,amsthm}
\usepackage{graphicx}
\usepackage[hidelinks]{hyperref}

\newtheorem{theorem}{Theorem}
\newtheorem{lemma}{Lemma}
\newtheorem{corollary}{Corollary}
\newcommand{\C}{\mathbb C}
\newcommand{\M}{M}
\newcommand{\norm}[1]{\left\lVert #1\right\rVert}

\title{A Flat-Discrepancy Prefix Criterion for the\\
Commutator Majorization Conjecture}
\author{Run \texttt{fa\_banach\_001}, lane 7}
\date{August 9, 2026}

\begin{document}
\maketitle

\section*{Status}

This packet records a likely valid partial result for Conjecture 7.1 of
Habib Zadeh--Sra \cite{HZS}.  It proves the conjectured inequality at any
prefix where one input has constant discrepancy values.  In particular it
settles the first two prefixes when one input is a normal \(3\times3\) matrix.
It does not settle the final trace-norm prefix in that case and is not a full
solution of the conjecture.

\section*{Source problem}

For \(T\in\M_n(\C)\), write

\[
  \norm{T}_{(k)}=\sum_{j=1}^k\sigma_j(T)
\]

for the Ky Fan \(k\)-norm.  The discrepancy values of \cite{HZS} are
characterized by

\begin{equation}
 \sum_{j=1}^k\delta_j(T)
   =\min_{z\in\C}\norm{T-zI}_{(k)}.
 \label{eq:variational}
\end{equation}

Conjecture 7.1 on source PDF page 20 asks whether, for all \(A,B\in\M_n(\C)\),

\begin{equation}
 \sigma([A,B])\prec_w 2\,\delta(A)\delta(B),
 \label{eq:conjecture}
\end{equation}

where the product on the right is entrywise.  Equivalently, for every
\(1\leq k\leq n\),

\[
 \norm{[A,B]}_{(k)}
 \leq 2\sum_{j=1}^k\delta_j(A)\delta_j(B).
\]

\begin{center}
\includegraphics[width=\textwidth]{figures/open_problem_crop.png}
\end{center}

\section*{Proof intuition}

The conjecture is invariant under independent scalar shifts of \(A\) and
\(B\).  At a fixed prefix \(k\), shift one matrix to minimize its operator
norm and the other to minimize its Ky Fan \(k\)-norm.  The ideal property of
Ky Fan norms then bounds the commutator by twice the product of those two
minima.  Formula \eqref{eq:variational} identifies the minima as
\(\delta_1(A)\) and \(\sum_{j\leq k}\delta_j(B)\).  If the first \(k\)
discrepancy values of \(A\) are equal, this product is exactly the conjectured
right-hand side, with no loss.

For a normal \(3\times3\) matrix the first two discrepancy values are
automatically equal.  This is a three-point planar fact: the minimum of the
largest eigenvalue distance is the smallest-enclosing-circle radius \(R\),
and the minimum of the sum of the two largest distances is \(2R\).

\section*{Formal result and proof}

We use only the standard Ky Fan ideal inequalities

\begin{equation}
 \norm{XY}_{(k)}\leq \norm{X}\norm{Y}_{(k)},\qquad
 \norm{YX}_{(k)}\leq \norm{Y}_{(k)}\norm{X}.
 \label{eq:ideal}
\end{equation}

Indeed, each follows from
\(\sigma_j(XY)\leq\norm{X}\sigma_j(Y)\), respectively its right-multiplication
version, after summing over \(j\leq k\).

\begin{theorem}[Flat-prefix criterion]
\label{thm:flat}
Let \(A,B\in\M_n(\C)\) and \(1\leq k\leq n\).  If

\[
  \delta_1(A)=\delta_2(A)=\cdots=\delta_k(A),
\]

then the \(k\)-th weak-majorization inequality in
\eqref{eq:conjecture} holds:

\[
 \norm{[A,B]}_{(k)}
 \leq 2\sum_{j=1}^k\delta_j(A)\delta_j(B).
\]

The same conclusion holds if the flat-prefix assumption is imposed on \(B\)
instead.
\end{theorem}

\begin{proof}
Put \(a=\delta_1(A)\).  By \eqref{eq:variational}, choose scalars
\(\alpha,\beta\in\C\) such that

\[
 \norm{A-\alpha I}=a,
 \qquad
 \norm{B-\beta I}_{(k)}=\sum_{j=1}^k\delta_j(B).
\]

The minima exist because the displayed norm functions are continuous and
coercive.  Set \(X=A-\alpha I\) and \(Y=B-\beta I\).  Scalar shifts do not
change a commutator, so the triangle inequality and \eqref{eq:ideal} give

\begin{align*}
 \norm{[A,B]}_{(k)}
 &=\norm{XY-YX}_{(k)}\\
 &\leq \norm{XY}_{(k)}+\norm{YX}_{(k)}\\
 &\leq 2\norm{X}\norm{Y}_{(k)}\\
 &=2a\sum_{j=1}^k\delta_j(B).
\end{align*}

Under the flat-prefix hypothesis, \(a=\delta_j(A)\) for every \(j\leq k\),
and the final expression is
\(2\sum_{j=1}^k\delta_j(A)\delta_j(B)\), as required.  Interchanging \(A\)
and \(B\) proves the symmetric assertion.
\end{proof}

\begin{lemma}[Three normal eigenvalues]
\label{lem:normal}
If \(A\in\M_3(\C)\) is normal and \(R(A)\) is the radius of the smallest
closed disk containing its spectrum, then

\[
 \delta_1(A)=\delta_2(A)=R(A).
\]
\end{lemma}

\begin{proof}
Let the eigenvalues be \(\lambda_1,\lambda_2,\lambda_3\).  Normality implies
that the singular values of \(A-zI\) are the three distances

\[
 r_i(z)=|\lambda_i-z|.
\]

Arrange these at each \(z\) as \(r_1\geq r_2\geq r_3\).  Formula
\eqref{eq:variational} for \(k=1\) gives

\[
 \delta_1(A)=\min_z r_1(z)=R(A).
\]

We claim that \(\min_z(r_1(z)+r_2(z))=2R(A)\).  For the lower bound, fix
\(z\), choose a farthest eigenvalue \(p\) with \(|p-z|=r_1\), and suppose
first that \(r_1>0\).  The disk centered at

\[
 c=z+\frac{r_1-r_2}{2r_1}(p-z)
\]

with radius \((r_1+r_2)/2\) contains \(p\).  It also contains each of the
other two eigenvalues \(q\), since

\[
 |q-c|\leq |q-z|+|z-c|
 \leq r_2+\frac{r_1-r_2}{2}
 =\frac{r_1+r_2}{2}.
\]

Thus \(R(A)\leq(r_1+r_2)/2\).  The case \(r_1=0\) is immediate.  Conversely,
a smallest disk containing a finite set has at least two boundary points
(with the zero-radius case again immediate), so at its center the two largest
distances both equal \(R(A)\).  This proves the claim.  Applying
\eqref{eq:variational} for \(k=2\) now yields

\[
 \delta_1(A)+\delta_2(A)=2R(A),
\]

and hence \(\delta_2(A)=R(A)=\delta_1(A)\).
\end{proof}

\begin{corollary}[Normal \(3\times3\) first two prefixes]
\label{cor:normal}
Let \(A\in\M_3(\C)\) be normal and let \(B\in\M_3(\C)\) be arbitrary.  Then
Conjecture 7.1 holds at \(k=1\) and \(k=2\):

\[
 \sum_{j=1}^k\sigma_j([A,B])
 \leq 2\sum_{j=1}^k\delta_j(A)\delta_j(B),
 \qquad k=1,2.
\]
\end{corollary}

\begin{proof}
The case \(k=1\) is Theorem \ref{thm:flat} with a one-term prefix.  Lemma
\ref{lem:normal} supplies the flat-prefix hypothesis for \(k=2\), so the same
theorem applies again.
\end{proof}

\section*{Verification and stress testing}

The formal proof reduces to three directly checkable ingredients:

\begin{enumerate}
 \item the scalar-shift variational formula \eqref{eq:variational}, which is
       Theorem 4.3 of the source paper;
 \item the Ky Fan triangle and ideal inequalities;
 \item the explicit enclosing-disk construction in Lemma \ref{lem:normal}.
\end{enumerate}

The included script \texttt{code/random\_search.py} independently minimizes
the scalar-shifted Ky Fan norms numerically and tests all majorization
prefixes.  Across 14,620 matrix pairs in dimensions \(3,4,5\), including
integer real and complex matrices, rank-one and rank-two matrices, triangular
and weighted-shift matrices, normal matrices, and perturbations of the sharp
pair \(E_{12},E_{21}\), it found no positive margin.  These experiments are
only a stress test; nonsmooth numerical minimization is not used in the proof
and does not establish the still-open prefixes.

\section*{Limitations and novelty check}

The result is prefixwise.  It proves no assertion at a prefix where neither
matrix has flat discrepancy values.  In particular, for a normal
\(3\times3\) matrix the third discrepancy value is generally controlled by a
geometric-median problem rather than the smallest enclosing circle.  The
argument therefore does not prove the \(k=3\) trace-norm inequality, even when
one input is normal, and the full conjecture remains open.

A bounded check on August 9, 2026 searched the run indexes, arXiv, and general
web indexes using arXiv:2111.11855, the exact title and authors, “Conjecture
7.1,” “discrepancy values,” “commutator majorization,” and normal
\(3\times3\) variants.  It found the source paper but no later full resolution
and no statement of the flat-prefix criterion or Corollary
\ref{cor:normal}.  Novelty confidence is moderate rather than definitive;
the claim should be reviewed as an elementary new consequence of the source
variational formula and standard Ky Fan inequalities.

\enlargethispage{5\baselineskip}
\paragraph{Human-review recommendation.}

Review as a likely-valid partial result.  The most useful checks are that the
two scalar shifts may indeed be chosen independently without changing the
commutator, and that the disk in Lemma \ref{lem:normal} covers all three
eigenvalues.  Do not promote this packet as a full solution unless the
remaining prefixes are supplied by a separate argument.

\begin{thebibliography}{9}

\bibitem{HZS}
P. Habib Zadeh and S. Sra,
\emph{Introducing discrepancy values of matrices with application to bounding
norms of commutators},
Linear Algebra Appl. 651 (2022), 359--384;
arXiv:2111.11855.

\end{thebibliography}

\end{document}
